\chapter{CLI, SDK, and API Quick Reference}
\index{Databricks CLI}\index{REST API}\index{quick reference}\index{SQL!reference}\index{PySpark!reference}%
This appendix collects the commands, endpoints, and idioms the book's labs use repeatedly, in
one lookup-friendly place. It is a reference, not a tutorial: each section points back to the
chapter that teaches the concept.

\section{Databricks CLI Essentials}
\index{Databricks CLI!essentials}%
\begin{lstlisting}[language=bash,caption={Databricks CLI commands used across the labs.},label={lst:appendixb-cli}]
databricks auth login --host https://<workspace>.azuredatabricks.net  # OAuth (Ch. 0)
databricks auth profiles                      # list configured profiles
databricks current-user me                    # confirm identity

# Workspace and clusters (Ch. 1)
databricks workspace list /
databricks clusters list --output json
databricks clusters create --json @cluster.json

# Jobs and pipelines (Ch. 12)
databricks jobs list
databricks jobs run-now <job-id>
databricks jobs submit --json @run.json
databricks runs repair --run-id <run-id>      # repair failed tasks only

# Unity Catalog (Ch. 8, 22)
databricks catalogs list
databricks schemas create bronze --catalog de_training
databricks volumes create de_training bronze landing MANAGED
databricks grants get schema de_training.gold

# Secrets (Ch. 21)
databricks secrets create-scope kv-lab
databricks secrets put-secret kv-lab kafka-password --string-value '<from-vault>'

# Asset Bundles (Ch. 23)
databricks bundle validate -t dev
databricks bundle deploy -t staging
databricks bundle run nightly_elt -t prod
\end{lstlisting}

\section{Databricks SDK for Python}
\index{databricks-sdk-py}%
The SDK mirrors the REST API and powers automation scripts throughout the book
(Chapters~1, 12, 23):

\begin{lstlisting}[language=Python,caption={SDK essentials.},label={lst:appendixb-sdk}]
from databricks.sdk import WorkspaceClient

w = WorkspaceClient()  # reads DATABRICKS_HOST / profile from env

me = w.current_user.me()
print(me.user_name)

# Submit a one-off job run (Ch. 1 smoke test)
from databricks.sdk.service.jobs import SubmitTask, NotebookTask
run = w.jobs.submit(run_name="ch0-smoke", tasks=[
    SubmitTask(task_key="smoke",
               notebook_task=NotebookTask(notebook_path="/Repos/<you>/course/src/smoke"),
               new_cluster={"spark_version": "15.4.x-scala2.12",
                            "num_workers": 1,
                            "node_type_id": "Standard_DS3_v2"})
]).result()

# List catalog tables (Ch. 8)
for t in w.tables.list(catalog_name="de_training", schema_name="gold"):
    print(t.full_name)
\end{lstlisting}

\section{REST API}
\index{REST API!Databricks}%
The REST API is the operations plane beneath Week~23's CI/CD. Pattern: acquire a token
(OAuth preferred), call the API, poll long-running operations.

\begin{lstlisting}[language=bash,caption={REST API essentials (token redacted).},label={lst:appendixb-rest}]
# List jobs
curl -H "Authorization: Bearer ******" `
    https://<workspace>.azuredatabricks.net/api/2.1/jobs/list

# Trigger a job run
curl -X POST -H "Authorization: Bearer ******" `
    -H "Content-Type: application/json" `
    -d '{"job_id": 12345}' `
    https://<workspace>.azuredatabricks.net/api/2.1/jobs/run-now

# Execute SQL on a warehouse
curl -X POST -H "Authorization: Bearer ******" `
    -H "Content-Type: application/json" `
    -d '{"statement":"SELECT COUNT(*) FROM de_training.gold.fact_orders","warehouse_id":"<id>"}' `
    https://<workspace>.azuredatabricks.net/api/2.0/sql/statements
\end{lstlisting}

\begin{pitfall}
The REST API authenticates as \emph{someone}: use a service principal with least-privilege
grants, an OAuth or vaulted secret, and a rotation date. Tokens pasted into scripts reproduce
Chapter~21's token-in-notebook failure at the platform layer.
\end{pitfall}

\section{Spark SQL and Delta Quick Reference}
\index{SQL!quick reference}\index{Delta Lake!quick reference}%
\begin{lstlisting}[language=SQL,caption={SQL idioms used across Chapters 2--8.},label={lst:appendixb-sql}]
-- Three-level namespace (Ch. 1, 8)
SELECT * FROM de_training.silver.orders LIMIT 10;

-- History and time travel (Ch. 2)
DESCRIBE HISTORY de_training.silver.orders;
SELECT COUNT(*) FROM de_training.silver.orders VERSION AS OF 12;
RESTORE TABLE de_training.silver.orders TO VERSION AS OF 11;

-- Maintenance (Ch. 5)
OPTIMIZE de_training.silver.orders ZORDER BY (customer_id, order_date);
ALTER TABLE de_training.gold.fact_orders CLUSTER BY (order_date, region);
VACUUM de_training.silver.orders RETAIN 168 HOURS;

-- Governance (Ch. 8, 22)
GRANT SELECT ON SCHEMA de_training.gold TO `analysts`;
CREATE OR REPLACE VIEW de_training.gold.v_customer_masked AS
SELECT customer_id,
       CASE WHEN is_member('pii_readers') THEN email
            ELSE sha2(email, 256) END AS email
FROM de_training.gold.dim_customer;

-- Cost introspection (Ch. 6, 23)
SELECT sku_name, SUM(usage_quantity) AS dbus
FROM system.billing.usage
WHERE usage_date >= current_date() - 30
GROUP BY sku_name ORDER BY dbus DESC;
\end{lstlisting}

\section{PySpark and Streaming Quick Reference}
\index{PySpark!quick reference}\index{Structured Streaming!quick reference}%
\begin{lstlisting}[language=Python,caption={PySpark/Streaming idioms used throughout.},label={lst:appendixb-pyspark}]
# Read a managed table; stamp lineage columns (Ch. 3)
df = (spark.read.table("de_training.bronze.orders")
      .withColumn("_ingested_at", F.current_timestamp()))

# Idempotent partition-scoped write (Ch. 9, 12, 20)
(df.write.format("delta").mode("overwrite")
   .option("replaceWhere", "as_of_date = '2026-07-31'")
   .saveAsTable("de_training.gold.features_daily"))

# Schema evolution (Ch. 2, 11): default enforcement rejects; opt in deliberately
(df.write.format("delta").mode("append")
   .option("mergeSchema", "true")
   .saveAsTable("de_training.silver.orders"))

# Auto Loader (Ch. 11)
stream = (spark.readStream.format("cloudFiles")
    .option("cloudFiles.format", "csv")
    .option("cloudFiles.schemaLocation",
            "/Volumes/de_training/bronze/_schemas/orders")
    .load("/Volumes/de_training/bronze/landing/orders"))

# Checkpointed streaming sink (Ch. 13)
(stream.writeStream.format("delta")
    .option("checkpointLocation", "/Volumes/de_training/bronze/_ckpt/orders")
    .trigger(availableNow=True)
    .toTable("de_training.bronze.orders_stream"))

# Watermark + dedup (Ch. 14)
deduped = (df.withWatermark("event_ts", "10 minutes")
             .dropDuplicates(["event_id"]))

# Registered model as UDF (Ch. 17, 20)
predict = mlflow.pyfunc.spark_udf(
    spark, "models:/ml.churn_model@champion", result_type="double")
\end{lstlisting}

\section{Asset Bundle Skeleton}
\index{Asset Bundles!quick reference}%
\begin{lstlisting}[language=json,caption={Minimal databricks.yml with dev/prod targets (Ch. 23).},label={lst:appendixb-bundle}]
bundle:
  name: governed-platform

variables:
  catalog:
    default: de_training

targets:
  dev:
    mode: development
    default: true
    workspace:
      host: https://<dev-workspace>.azuredatabricks.net
  prod:
    mode: production
    workspace:
      host: https://<prod-workspace>.azuredatabricks.net

resources:
  jobs:
    nightly_elt:
      name: nightly-elt-${bundle.target}
      tasks:
        - task_key: pipeline
          pipeline_task:
            pipeline_id: ${resources.pipelines.medallion.id}
  pipelines:
    medallion:
      name: medallion-${bundle.target}
      catalog: ${var.catalog}
      libraries:
        - notebook:
            path: ../src/medallion_pipeline.py
\end{lstlisting}

\begin{tipbox}
Keep this appendix open during capstones. Every idiom here appears in at least one chapter's
validation checklist; when an acceptance criterion says ``idempotent'' or ``dedup,'' these are the
exact patterns it refers to.
\end{tipbox}
